\documentclass[journal]{IEEEtran}
\usepackage{amsmath,amsfonts}
\usepackage{algorithmic}
\usepackage{algorithm}
\usepackage{array}
\usepackage[caption=false,font=normalsize,labelfont=sf,textfont=sf]{subfig}
\usepackage{textcomp}
\usepackage{stfloats}
\usepackage{url}
\usepackage{verbatim}
\usepackage{graphicx}
\usepackage{cite}
% ===== Preamble (put these once) =====
\usepackage{booktabs,multirow,colortbl,xcolor}
\usepackage{graphicx}   % for \rotatebox
\usepackage{makecell}   % for stacked headers
\usepackage{amssymb}    % for \checkmark
\usepackage{arydshln}
\usepackage[table]{xcolor} % enables \cellcolor in tables
\usepackage{adjustbox} 
% Define nice light tones (you can tweak these)
\definecolor{lightgreen}{HTML}{DFF0D8} % pastel green
\definecolor{lightyellow}{HTML}{FCF8E3} % pastel yellow

\definecolor{errbg}{RGB}{255,230,230} % light red
\definecolor{accbg}{RGB}{230,245,235} % light green
\definecolor{probg}{RGB}{233,238,251} % light blue
\providecommand{\best}[1]{\textbf{#1}} % bold-best helper used in result tables
\hyphenation{op-tical net-works semi-conduc-tor IEEE-Xplore}
% updated with editorial comments 8/9/2021

\begin{document}

\title{Self-Supervised Metric Adaptation of Depth Foundation Models via Scale--Shift Decoupling}

\author{%
    Fardin Ayar, Department of Computer Science, Amirkabir University of Technology\\
    \and
    Hossein Vahidi, Ferdowsi University of Mashhad\\
    \and
    Najmeh Mohammadbagheri, Amirkabir University of Technology\\
    \and
    Mahdi Javanmardi*, Amirkabir University of Technology\\
    \and
    Francesco Pilla, University College Dublin\\
}



%\author{%
%	Behnam~Roshanfekr,
%	Mohammad~Rahmati,
%	Maryam~Amirmazlaghani, and
%	Maryam~Abdolali%
%	\thanks{B. Roshanfekr, M. Rahmati (corresponding author), and M. Amirmazlaghani are with the Department of Computer Engineering, Amirkabir University of Technology (Tehran Polytechnic), Iran (e-mails: b.roshanfekr@aut.ac.ir; rahmati@aut.ac.ir; mazlaghani@aut.ac.ir).}%
%	\thanks{M. Abdolali is with the Department of Mathematics and Computer Science, K.N.Toosi University (KNTU), Iran (e-mail: maryam.abdolali@kntu.ac.ir).}%
%}

%\author{Behnam Roshanfekr,~\IEEEmembership{Staff,~IEEE,}
        % <-this % stops a space
%\thanks{This paper was produced by the IEEE Publication Technology Group. They are in Piscataway, NJ.}% <-this % stops a space
%\thanks{Manuscript received April 19, 2021; revised August 16, 2021.}}

% The paper headers
\markboth{Journal of \LaTeX\ Class Files,~Vol.~14, No.~8, August~2021}%
{Shell \MakeLowercase{\textit{et al.}}: A Sample Article Using IEEEtran.cls for IEEE Journals}

%\IEEEpubid{0000--0000/00\$00.00~\copyright~2021 IEEE}
% Remember, if you use this, you must call \IEEEpubidadjcol in the second
% column for its text to clear the IEEEpubid mark.

\maketitle

\begin{abstract}
    Affine-invariant depth foundation models predict relative depth with remarkable zero-shot generalization and exceptional edge and fine-structure fidelity, but their predictions lack absolute (metric) scale, which is essential for navigation and obstacle avoidance. Recovering metric scale typically demands LiDAR or ground-truth depth, or switching to metric foundation models that sacrifice geometric detail. We instead ask whether a frozen affine-invariant foundation model can be turned \emph{metric} purely through self-supervision, without any depth annotations. The central obstacle is a fundamental tension: self-supervised photometric losses are scale-invariant, so naive fine-tuning not only fails to learn absolute scale but actively erodes the rich relative structure of the pretrained model by collapsing it toward scale-ambiguous solutions. Our key idea is to \emph{decouple metric scale from depth shape}: the frozen decoder keeps producing the foundation model's sharp relative prediction, while a lightweight per-image affine transformation---regressed from the encoder's class tokens---carries the global metric scale and shift. Absolute scale is injected only through the pose branch via readily available GPS traces, and a teacher distillation loss from the frozen foundation model protects fine structure from photometric erosion. The entire adaptation is realized with Low-Rank Adaptation, training under $1\%$ of parameters. On KITTI, our single-image model produces metric depth that matches or surpasses DepthAnything\,V2 with substantially larger backbones, approaches a supervised metric foundation model, and rivals recent self-supervised methods, while preserving strong cross-domain transfer and remaining effective when trained on as little as $1\%$ of the data.
\end{abstract}

\begin{IEEEkeywords}
Monocular depth estimation, self-supervised learning, depth foundation models, metric depth, parameter-efficient fine-tuning, scale recovery.
\end{IEEEkeywords}

\section{Introduction}\label{sec1}

Monocular depth estimation is a core component in robotics, autonomous navigation, and 3D scene understanding. Self-Supervised Monocular Depth (SMD) methods~\cite{zhou2017unsupervised,8100182,godard2019digging} remove the need for costly ground-truth depth labels by exploiting photometric consistency between adjacent video frames, training depth and pose networks with view-synthesis losses under geometric warping. Despite steady progress, SMD methods inherit two persistent weaknesses. First, the photometric objective is fundamentally ambiguous with respect to scale: depth and ego-motion can be multiplied by any common factor without changing the reprojection error, so absolute (metric) depth cannot be recovered. Second, because these models are trained from scratch on a single dataset, they learn depth representations from a narrow distribution and generalize poorly across domains, often failing under shifts in camera, scene statistics, or appearance.

Concurrently, depth foundation models such as MiDaS~\cite{ranftl2020towards}, Metric3D~\cite{liu2024metricdepth}, and DepthAnything\,V2~\cite{yang2024depth} have demonstrated striking zero-shot generalization by training on massive, heterogeneous data. The most accurate of these are \emph{affine-invariant}: trained with scale- and shift-invariant losses, they predict \emph{relative} depth and, freed from the burden of committing to an absolute scale, capture fine-grained geometry, edges, and thin structures with unmatched fidelity. This very property is also their main practical limitation: their predictions lack metric scale, which navigation and obstacle avoidance require. Metric foundation models~\cite{bhat2023zoedepth,yin2023metric3d,piccinelli2024unidepth} predict absolute depth directly, but enforcing a metric constraint during training measurably degrades edge sharpness and fine detail relative to their affine-invariant counterparts~\cite{li2025benchdepth}. Closing this gap by converting an affine-invariant model to metric typically reintroduces exactly the data dependency foundation models were meant to remove---LiDAR or ground-truth depth for scale calibration.

This exposes a compelling opportunity. Affine-invariant foundation models provide the best relative geometry and the strongest cross-domain priors, yet lack the one ingredient---metric scale---that downstream systems need, while self-supervision offers a label-free route to scale through camera geometry. We therefore ask: \emph{can a frozen affine-invariant depth foundation model be adapted to metric, single-image depth purely through self-supervision, without any depth annotations and without sacrificing the geometric fidelity that makes it valuable?}

Answering this question is harder than it appears, because of a fundamental tension that we place at the center of this work. Self-supervised photometric losses are scale-invariant, so naively fine-tuning a pretrained model to minimize reprojection error does not merely fail to learn absolute scale---it actively \emph{erodes} the metric and relative structure of the foundation model by driving its parameters toward scale-ambiguous solutions. Any successful adaptation must simultaneously (i) introduce a route for absolute scale that the photometric loss alone cannot provide, and (ii) protect the pretrained geometry from being washed out by that same loss.

We resolve this tension with a single guiding principle---\emph{decouple metric scale from depth shape}---instantiated by three lightweight, complementary components. \textbf{First}, a scale--shift decoupling head lets the frozen decoder keep emitting the foundation model's sharp relative prediction, while a per-image affine transformation, regressed from the encoder's class (\texttt{CLS}) tokens, carries the two global degrees of freedom---scale and shift---needed to make it metric. Because the affine acts on a standardized depth field, gradients adjust \emph{scale} without distorting \emph{shape}. \textbf{Second}, absolute scale is injected only through the pose branch, by aligning predicted ego-motion magnitude with displacement from readily available GPS traces~\cite{chawla2021multimodal}; this metric reference then propagates to depth through the shared warping geometry. \textbf{Third}, a teacher distillation loss from the frozen foundation model, applied in a patch-wise scale-shift-invariant form, explicitly protects edges and thin structures from photometric erosion---and, being itself scale-shift-invariant, constrains shape without fighting the scale signal. A curriculum couples the two so distillation anchors structure early while GPS scale ramps in. The full adaptation is realized with Low-Rank Adaptation (LoRA)~\cite{hu2022lora}, training under $1\%$ of parameters and thereby preserving the model's zero-shot generalization. Our contributions are:

\begin{itemize}
    \item \textbf{Self-supervised metric adaptation of a relative foundation model.} We present, to our knowledge, the first framework that turns a frozen affine-invariant depth foundation model into a \emph{metric}, single-image predictor using only self-supervision and GPS---no depth ground truth, no LiDAR---while retaining its cross-domain generalization.
    \item \textbf{Scale--shift decoupling.} We introduce a \texttt{CLS}-conditioned per-image affine head that separates the global metric component (scale and shift) from the dense relative geometry, so metric supervision can act on a low-dimensional, well-conditioned target without degrading the foundation model's fine structure.
    \item \textbf{A training recipe that resolves the scale-invariance/erosion tension.} We combine GPS-anchored scale supervision, teacher scale-shift-invariant distillation, and a curriculum that couples them, together with a robust high-loss masking scheme that suppresses dynamic objects and occlusions without optical flow or semantics.
    \item \textbf{Strong, data-efficient results.} Our single-image model matches or surpasses DepthAnything\,V2 with larger backbones, approaches a supervised metric foundation model, rivals recent self-supervised methods, generalizes across domains, and remains effective with as little as $1\%$ of the training data.
\end{itemize}


\section{Related work}\label{sec2}
\subsection{Self-Supervised Monocular Depth Estimation}

Motivated by the prohibitive annotation cost of depth datasets and inspired by stereo-supervised approaches~\cite{garg2016unsupervised, 8100182}, self-supervised monocular depth estimation (SMD) has emerged as a compelling alternative to supervised learning. The seminal work of Zhou et al.~\cite{zhou2017unsupervised} established the foundational framework by jointly training depth and pose networks to minimize a photometric view-synthesis loss, reconstructing each target frame from adjacent source frames under geometric warping. This formulation relies on two assumptions---scene staticity and the absence of occlusions---that are frequently violated in practice, producing artifacts around moving objects and disoccluded regions. Subsequent work introduced a range of mitigations: minimum-reprojection and auto-masking heuristics to discard unreliable pixels~\cite{godard2019digging}, and auxiliary signals such as optical flow~\cite{casser2019depth, chen2019self, hui2022rm, nguyen2024mining, sun2024dynamo}, scene flow~\cite{casser2019depth}, and semantic segmentation~\cite{klingner2020self} to reason about dynamic motion explicitly. Architectural advances---depth-aware convolutions~\cite{arch1, arch2} and transformer backbones~\cite{monovit, arch3, wang2024sqldepth}---further improved accuracy by modeling long-range context. A number of methods additionally aggregate several frames at test time, typically through cost volumes that encode cross-frame correspondence~\cite{watson2021temporal, woo2024prodepth}; while accurate, these models add computational cost and, like all SMD methods, remain bound by two structural limitations that motivate our work: the photometric objective is scale-ambiguous, so it cannot recover metric depth, and networks trained from scratch on a single dataset generalize poorly beyond their training distribution.

\subsection{Depth Foundation Models}

The success of large-scale pretraining in vision and language has inspired depth foundation models that learn generalizable depth priors from massive, heterogeneous data. Pioneering works such as MegaDepth~\cite{li2018megadepth} and MiDaS~\cite{ranftl2020towards} showed that training on diverse image collections with scale- and shift-invariant (SSI) losses yields impressive zero-shot prediction across unseen domains, and DepthAnything and its successor DepthAnything\,V2~\cite{yang2024depth} pushed affine-invariant estimation to state-of-the-art fidelity in capturing edges and fine geometric detail. Crucially, the SSI training objective is what enables this fidelity: by never committing to an absolute scale, the model is free to allocate capacity to relative structure. The same objective, however, leaves predictions without metric scale, limiting their direct use in navigation and manipulation. Metric foundation models address this by predicting absolute depth---ZoeDepth~\cite{bhat2023zoedepth} via relative-to-metric heads, Metric3D~\cite{yin2023metric3d} and its successor~\cite{liu2024metricdepth} by conditioning on camera intrinsics, and UniDepth~\cite{piccinelli2024unidepth} by pursuing universal metric prediction---but recent benchmarking~\cite{li2025benchdepth} reports that enforcing a metric constraint during training consistently degrades geometric detail relative to affine-invariant models, revealing a trade-off between metric consistency and representation richness. Our work aims to obtain the best of both: the geometric fidelity of an affine-invariant model and the metric scale of a metric one.

\subsection{Parameter-Efficient Adaptation and Metric Scale Recovery}

Rather than training a metric model from scratch, we adapt a frozen affine-invariant foundation model, which requires two ingredients: an adaptation mechanism that does not destroy the pretrained priors, and a source of metric scale that does not rely on depth annotations. For the former, parameter-efficient fine-tuning---most prominently Low-Rank Adaptation (LoRA)~\cite{hu2022lora}---updates a small fraction of weights through low-rank residuals, enabling task adaptation while mitigating catastrophic forgetting; we use it to graft self-supervised, metric behavior onto the foundation model while keeping its backbone frozen. For the latter, while metric supervision is classically obtained from LiDAR or ground-truth depth, GPS and other ego-motion sensors provide an inexpensive, widely available metric reference: Chawla et al.~\cite{chawla2021multimodal} align predicted translation magnitude with GPS-derived displacement to resolve the scale ambiguity of self-supervised pose. We build on this idea but apply it in the setting of foundation-model adaptation, and pair it with an explicit \emph{scale--shift decoupling} mechanism so that this weak metric signal can set the global scale of a relative prediction without disturbing its shape---directly targeting the tension between metric scale and geometric fidelity that prior approaches leave unaddressed.

Closest to our goal of annotation-free metric monocular depth for driving are MonoPP~\cite{elazab2025monopp} and its successor Gamma-from-Mono~\cite{elazab2025gamma}, which recover metric scale in closed form from planar-parallax geometry and the camera's known mounting height above the road plane---MonoPP via a multi-frame teacher that distills static-scene scale into a single-image student, Gamma-from-Mono via a single compact network that regresses a dimensionless plane-deviation field (\(\gamma\)) and converts it to metric depth deterministically. Both, however, train lightweight networks from scratch and hinge on a strong, explicit geometric assumption: a visible, locally planar road surface and a known, calibrated camera pose relative to it. Our route to metric scale is comparatively weak---a single GPS-derived displacement ratio with no scene-geometry assumptions---but this is precisely what lets it generalize beyond flat, well-characterized road scenes (we additionally demonstrate zero-shot transfer to Cityscapes, Section~4.5), and what lets us graft it onto a large pretrained backbone instead of learning depth representations from scratch, inheriting the foundation model's edge fidelity and cross-domain transfer that compact from-scratch networks do not exhibit.

Mechanistically, the work closest to our scale--shift decoupling head is RSA~\cite{zeng2024rsa}, which likewise recovers metric depth by regressing a single global affine transformation that is applied on top of a \emph{frozen} relative-depth backbone---the same separation of ``what carries scale'' from ``what carries shape'' that we exploit. We find this convergence telling rather than discouraging: it suggests the decoupling principle is the right one for adapting frozen relative predictors. The two methods differ, however, in exactly the places that matter for an annotation-free pipeline. RSA conditions its affine on a language caption encoded by a frozen CLIP model and, more importantly, \emph{trains} the affine regressor with a supervised L1 loss against ground-truth metric depth---reintroducing the very annotation dependency that motivates working with foundation models in the first place. Our affine is instead conditioned on the depth encoder's own class tokens (no auxiliary modality or caption pipeline) and is learned purely from photometric reprojection and GPS, with no metric labels at any stage---an honest but, we believe, important difference: the same correction mechanism, but one that can be deployed wherever a camera and a coarse displacement signal are available, and the other that still requires a calibration set with ground-truth depth.


\section{Proposed method}\label{sec3}
Our framework adapts an affine-invariant depth foundation model into a \emph{metric}, single-image predictor through lightweight modifications that preserve the pretrained representation. We build upon DepthAnything\,V2~\cite{yang2024depth}, keeping its encoder and decoder frozen and introducing four components: (i) parameter-efficient LoRA modules~\cite{hu2022lora} that adapt the frozen backbone without catastrophic forgetting, (ii) a scale--shift decoupling head that separates the global metric scale and shift---regressed from the encoder's class tokens---from the dense relative geometry produced by the decoder, (iii) a pose network that estimates inter-frame ego-motion and receives GPS-based metric supervision, and (iv) a frozen teacher copy of the foundation model that distills fine structure into the adapted model. Adjacent frames and GPS are used \emph{only at training time}; at inference the model predicts metric depth from a single image. We first review the self-supervised preliminaries, then describe each component and the training objective.

\subsection{Preliminaries: Self-Supervised Depth Estimation}

Following the self-supervised paradigm~\cite{zhou2017unsupervised, godard2019digging}, we exploit photometric consistency between temporally adjacent training frames as the primary supervisory signal. Given a target frame \(I_{t}\) and source frames \(\{I_{t+n}\mid n \in \{-1, 1\}\}\), we predict the target depth \(D_t\) and the relative camera pose \(T_{t \to t+n}\). These enable differentiable view synthesis~\cite{garg2016unsupervised}: we reconstruct the target from each source by
\begin{equation}
  \hat{I}_{t}(p) = I_{t+n}(p'), \quad \text{where } p' \sim K\, T_{t \to t+n}\, D_t(p)\, K^{-1} p,
\end{equation}
with \(K\) the camera intrinsics and \(p\) pixel coordinates. The photometric reconstruction loss penalizes the discrepancy between the synthesized and observed target:
\begin{equation}
    \mathcal{L}_{\text{photo}} = \sum_{p} \rho\big(I_t(p), \hat{I}_{t}(p)\big),
\end{equation}
where \(\rho(\cdot)\) is a robust photometric error combining structural similarity (SSIM) and \(\ell_1\) distance~\cite{godard2019digging}. We handle violations of the static-scene assumption with auto-masking and a robust high-loss rejection scheme (Section~3.6). Crucially, this objective is invariant to a common rescaling of depth and translation, a property we must confront directly to recover metric scale.

\bigbreak
\begin{figure*}
    \centering
    % NOTE: regenerate this figure to match the single-image pipeline described below.
    \includegraphics[width=\linewidth]{imgs/Picture1.png}
    \caption{Overview of the proposed self-supervised metric adaptation framework.
    A single target frame $I_t$ is encoded by a frozen DepthAnything\,V2 backbone (adapted with LoRA) and decoded by the DPT head into a relative disparity field. A \texttt{CLS}-conditioned scale--shift head predicts a per-image affine $(s,b)$ that converts the relative prediction into metric depth $D_t$, decoupling the global scale from the dense shape. A PoseNet estimates inter-frame ego-motion and is supervised by GPS displacement to inject metric scale, which propagates to depth through the photometric warping. A frozen teacher copy of the foundation model distills fine structure via a scale-shift-invariant loss. Adjacent frames, the pose branch, and the teacher are used only during training; inference is single-image.}
    \label{fig:placeholder}
\end{figure*}


\subsection{Architecture and the Scale-Invariance Tension}

\textbf{Foundation model selection.} While metric foundation models exist~\cite{yin2023metric3d,piccinelli2024unidepth}, we deliberately build on the \emph{affine-invariant} DepthAnything\,V2~\cite{yang2024depth} for its superior edge preservation and geometric detail~\cite{li2025benchdepth}, and accept the challenge of recovering metric scale through self-supervision. This choice is governed by a fundamental tension: self-supervised photometric losses are scale-invariant, so naive fine-tuning will not only fail to learn absolute scale but will actively erode the pretrained structure by driving parameters toward scale-ambiguous solutions. Our architecture is designed to address both halves of this tension---providing a controlled locus for metric scale while shielding the dense geometry from photometric erosion.

\textbf{Feature extraction.} Let \(E_{\mathbf w}\) denote the Vision Transformer encoder. For an input image \(I_{t} \in \mathbb{R}^{3 \times H \times W}\), we extract multi-scale features from four layers following the DPT architecture, together with the corresponding class tokens:
\begin{equation}
\big(\{F^{i}_t\}_{i=1}^{4},\, \{c^{i}_t\}_{i=1}^{4}\big) = E_{\mathbf w}(I_t),
\qquad F^{i}_t \in \mathbb{R}^{C \times \frac{H}{14} \times \frac{W}{14}},\; c^{i}_t \in \mathbb{R}^{C}.
\end{equation}
The spatial features are decoded by the DPT decoder \(D_{\boldsymbol\phi}\) into a relative disparity field \(\hat d_t = D_{\boldsymbol\phi}(\{F^{i}_t\}) \in \mathbb{R}^{1\times H\times W}\), while the class tokens \(\{c^{i}_t\}\) drive metric scale recovery (Section~3.3).

\textbf{Parameter-efficient adaptation via LoRA.} To preserve the zero-shot generalization of DepthAnything\,V2, we freeze the encoder \(E_{\mathbf w}\) and decoder \(D_{\boldsymbol\phi}\), and insert Low-Rank Adaptation (LoRA)~\cite{hu2022lora} modules at (i) the query and key projections of the encoder's self-attention layers and (ii) the convolutional layers of the DPT decoder. LoRA reparameterizes each adapted weight as \(\mathbf{W} + \mathbf{B}\mathbf{A}\) with \(\mathrm{rank}(\mathbf{B}\mathbf{A}) = r \ll d\), making fewer than \(1\%\) of parameters trainable. This constrains the adapted model to stay close to the pretrained one---a structural safeguard against the erosion described above---while still allowing it to learn metric behavior.

\subsection{Scale--Shift Decoupling}

The defining feature of an affine-invariant model is that its prediction is meaningful only up to an unknown global scale and shift; equivalently, exactly two global degrees of freedom separate its relative output from a metric one. Our central design principle is to make these two degrees of freedom \emph{explicit} and to route all metric supervision through them, leaving the high-dimensional shape untouched.

Let \(\hat d_t\) be the decoder's relative disparity. We standardize it using its own per-image statistics, treated as constants (stop-gradient, denoted \(\lfloor\cdot\rfloor\)):
\begin{equation}
\bar d_t = \frac{\hat d_t - \lfloor\mu_t\rfloor}{\lfloor\sigma_t\rfloor},
\qquad \mu_t = \mathrm{mean}(\hat d_t),\;\; \sigma_t = \mathrm{std}(\hat d_t).
\end{equation}
A lightweight head \(g_{\boldsymbol\theta}\) regresses a per-image affine from the concatenated class tokens, summarizing global scene context (illumination, layout, scale cues) without spatial bias:
\begin{equation}
(s_t,\, b_t) = g_{\boldsymbol\theta}\big([\,c^{1}_t; c^{2}_t; c^{3}_t; c^{4}_t\,]\big),
\qquad s_t = \mathrm{softplus}(\cdot) > 0,\;\; b_t \in \mathbb{R},
\end{equation}
and applies it to the standardized field to produce the corrected disparity, from which metric depth follows by inversion:
\begin{equation}
d_t = \lfloor\sigma_t\rfloor\,(s_t\, \bar d_t + b_t) + \lfloor\mu_t\rfloor,
\qquad
D_t = \big(\max(d_t,\, 1/d_{\max})\big)^{-1}.
\label{eq:scaleshift}
\end{equation}
Two properties make this an effective decoupling. \emph{First}, because \(\mu_t,\sigma_t\) are detached, the only path by which a global rescaling of the output can be realized is through the scalars \((s_t,b_t)\); a metric loss therefore adjusts scale and shift through two well-conditioned parameters rather than by redistributing magnitude across the dense field, so it cannot reshape the relative geometry to satisfy scale. \emph{Second}, conditioning on the class tokens lets the per-image affine adapt to scene content (e.g.\ open road vs.\ cluttered intersection) while keeping the dense decoder identical to the frozen foundation model. We zero-initialize the scale and shift heads so that training begins from the near-identity affine \((s_t\!\approx\!1, b_t\!\approx\!0)\), i.e.\ from the pretrained relative prediction, and departs from it only as much as the metric signal requires. The head adds a negligible number of parameters and no inference-time cost beyond a small MLP.

\subsection{GPS-Based Metric Scale Recovery}

Scale--shift decoupling provides a \emph{place} for metric scale to live; we now provide the \emph{signal}. Since the photometric loss is scale-invariant, absolute scale must enter from an external metric reference. Following~\cite{chawla2021multimodal}, we use GPS traces---ubiquitous in driving data---to supervise the magnitude of the pose network's predicted translation. GPS positions are converted to a local Cartesian frame (ENU/UTM) yielding \(\mathbf{p}_k \in \mathbb{R}^3\) in meters; for a source--target pair the displacement is \(\mathbf{G}_{s \to t} = \mathbf{p}_t - \mathbf{p}_s\), and the pose network predicts relative translation \(\hat{\mathbf{T}}_{s \to t}\). We align their magnitudes,
\begin{equation}
\mathcal{L}_{\mathrm{g2s}}
=
\sum_{s,t}
\Bigg(
\frac{\lVert \mathbf{G}_{s\!\rightarrow t}\rVert_2}
     {\lVert \hat{\mathbf{T}}_{s\!\rightarrow t}\rVert_2}
-1
\Bigg)^{\!2},
\label{eq:gps_scale_loss}
\end{equation}
constraining only the scale, not the direction, which can be noisy under GPS error or changes in heading. To make supervision robust we gate it on frames with sufficient measured motion---skipping near-stationary pairs where the GPS ratio is ill-conditioned---and apply the constraint in log space in practice. GPS acts \emph{solely on the pose network}; the metric scale it fixes then propagates to depth through the shared warping geometry of the photometric loss, which couples depth and translation scale. Concretely, once the translation magnitude is metric, minimizing reprojection error drives the depth scale---and hence the affine \(s_t\) of Eq.~\eqref{eq:scaleshift}---toward the consistent metric value. Unlike LiDAR supervision, GPS is lightweight, widely available, and sensor-free.

\subsection{Teacher-Guided Scale-Shift-Invariant Distillation}

Even with adaptation constrained by LoRA, prolonged photometric training tends to soften edges and thin structures---the very features for which we chose an affine-invariant backbone. We counteract this with a distillation loss from a \emph{frozen} copy of DepthAnything\,V2, which serves as a teacher of fine structure. Because the teacher is itself affine-invariant, a naive pixel-wise loss against it would inject an arbitrary scale and shift and fight the metric signal of Section~3.4. We therefore distill only \emph{shape}, in a patch-wise scale-shift-invariant (SSI) form.

For each \(K\times K\) patch \(p\) (\(K=32\), \(50\%\) overlap), we \(z\)-normalize the student and teacher disparities within the patch and compare them, making the term invariant to local affine transformations and tolerant of sub-pixel misalignment:
\begin{equation}
\ell^{\mathrm{SSI}}_{p}
= \frac{1}{K^{2}} \sum_{(u,v)\in p} \big|\,\bar d^{\mathrm{stu}}_{p}(u,v)-\bar d^{\mathrm{tea}}_{p}(u,v)\,\big|,
\end{equation}
where \(\bar{(\cdot)}_p\) denotes per-patch standardization. We evaluate the loss only on high-variance patches \(\mathcal{P}=\{\,p \mid \operatorname{Var}(p)>\tau\,\}\), concentrating supervision on depth discontinuities and thin structures rather than flat regions:
\begin{equation}
\mathcal{L}_{\mathrm{SSI}}
= \frac{1}{|\mathcal{P}|}\sum_{p\in\mathcal{P}} \ell^{\mathrm{SSI}}_{p}.
\end{equation}
An optional whole-image SSI term constrains global depth ordering and far-field structure that lie beyond a single patch. Being scale-shift-invariant, this distillation protects geometric fidelity while remaining orthogonal to metric scale: it constrains \emph{shape}, GPS constrains \emph{scale}, and the two do not interfere by construction---an explicit realization of our decoupling principle on the loss side.

\textbf{Curriculum coupling.} The two objectives have complementary roles over training: structure should be anchored early, while the pose network and scale head stabilize before metric scale is enforced. We therefore couple their schedules through a single curriculum. The GPS term is ramped \emph{up} with a weight \(w_{\mathrm{GPS}}(t)\) that grows over training, preventing premature scale constraints while the pose network is still noisy; the distillation weight is correspondingly ramped \emph{down}, \(\lambda_{\mathrm{SSI}}(t) \propto 1 - w_{\mathrm{GPS}}(t)\), so structural anchoring is strongest at initialization and yields as the model commits to a metric scale. This dynamic weighting ties the two halves of the tension into a single, self-consistent schedule.

\subsection{Robust Photometric Training}

\textbf{Auto-masking.} Following Monodepth2~\cite{godard2019digging}, the per-pixel photometric error combines SSIM and \(\ell_1\):
\begin{equation}
pe(I_a, I_b) = \frac{\alpha}{2}\,\big(1-\text{SSIM}(I_a, I_b)\big)
+ (1-\alpha)\,\lVert I_a - I_b \rVert_1,
\qquad \alpha = 0.85,
\label{eq:pe}
\end{equation}
and we discard pixels whose warped error exceeds the stationary (identity) error,
\begin{equation}
\mu = \Big[\, \min_{n} pe\!\big(I_t,\; \hat I_{t}^{\,n}\big) \;<\; \min_{n} pe\!\big(I_t,\; I_{t+n}\big) \,\Big],
\label{eq:auto-mask}
\end{equation}
with \([\cdot]\) the Iverson bracket and \(\hat I_{t}^{\,n}\) the reconstruction from source \(n\).

\textbf{High-loss rejection.} Auto-masking removes stationary and ego-aligned pixels but not dynamic objects moving with the camera, severe occlusions, or non-Lambertian regions, which produce large but non-trivial reprojection errors. Rather than relying on optical flow or semantic labels, we suppress these outliers directly: among all pixels that survive auto-masking, we rank reprojection errors across the batch and discard the top \(\eta\) fraction (the highest-error tail). This flow- and semantics-free scheme prevents a small number of hard, typically dynamic, pixels from dominating the gradient. The final masked reprojection loss takes the minimum over sources on the retained pixels:
\begin{equation}
\mathcal{L}_{\text{reproj}} = \frac{\sum_p m(p)\,\min_{n} pe\!\big(I_t(p),\; \hat I_{t}^{\,n}(p)\big)}{\sum_p m(p) + \epsilon},
\label{eq:Lreproj}
\end{equation}
where \(m\) is the combined auto-mask and high-loss-rejection mask.

\subsection{Training Objective and Optimization}

The total objective combines the photometric reconstruction loss with GPS metric supervision, teacher SSI distillation, and an edge-aware smoothness regularizer:
\begin{equation}
\mathcal{L}_{\text{total}} = \mathcal{L}_{\text{reproj}}
+ \lambda_{\mathrm{GPS}}(t)\,\mathcal{L}_{\mathrm{g2s}}
+ \lambda_{\mathrm{SSI}}(t)\,\mathcal{L}_{\mathrm{SSI}}
+ \lambda_{\mathrm{sm}}\,\mathcal{L}_{\mathrm{sm}},
\end{equation}
where \(\lambda_{\mathrm{GPS}}(t)\) and \(\lambda_{\mathrm{SSI}}(t)\) follow the coupled curriculum of Section~3.5. We train the LoRA parameters, the scale--shift head \(g_{\boldsymbol\theta}\), the final DPT decoder layers, and the pose network (initialized from a Cityscapes-pretrained pose encoder); the encoder and decoder backbones, as well as the teacher copy, remain frozen. This keeps trainable parameters below \(1\%\) of the model, which we find to be the key to preserving the foundation model's generalization while grafting metric, self-supervised behavior onto it.




\section{Experiments}\label{sec4}




We evaluate our framework along several axes: comparison with depth foundation models, comparison with state-of-the-art self-supervised methods, cross-domain transfer, data efficiency under limited training data, and a component ablation. Throughout, our model performs \emph{single-image} inference; adjacent frames and GPS are used only during training. Our experiments show that the proposed adaptation grafts metric scale onto an affine-invariant foundation model while preserving its geometric fidelity and generalization.

\subsection{Experimental Setup}

\textbf{Datasets.} We conduct experiments on two driving benchmarks, both featuring dynamic objects and diverse conditions.
\bigbreak

\textbf{KITTI}~\cite{geiger2013vision} is our primary benchmark, providing imagery and LiDAR depth from driving scenarios. We adopt the standard Eigen split~\cite{eigen2015predicting} for training and testing, and additionally evaluate on the improved KITTI ground truth~\cite{uhrig2017sparsity}, which provides denser annotations than the raw sparse LiDAR projections.

\textbf{Cityscapes}~\cite{cordts2016cityscapes} enables zero-shot cross-domain evaluation. It comprises urban driving sequences from 50 European cities with visual characteristics substantially different from KITTI. Following standard protocols~\cite{zhou2017unsupervised, gordon2019depth}, we evaluate on 1,525 test frames using SGM disparity maps~\cite{hirschmuller2007stereo} as ground truth, with depth capped at 80m.
\bigbreak

\begin{table*}[t]
\centering
\caption{Comparison with recent self-supervised methods on the KITTI Eigen split. ``Metric'' (\checkmark) denotes direct metric output; otherwise predictions are scale-aligned by median scaling. Best in \colorbox{lightgreen}{green}, second best in \colorbox{lightyellow}{yellow}.}
\label{tab:selfsupervised}
\setlength{\tabcolsep}{4pt}
\renewcommand{\arraystretch}{1.1}
\resizebox{\linewidth}{!}{%
\begin{tabular}{l|c|c|cccc|ccc}
\toprule
\multirow{2}{*}{\textbf{Method}} &
\multirow{2}{*}{\textbf{Metric}} &
\multirow{2}{*}{Res.} &
\multicolumn{4}{c|}{\cellcolor{errbg}\textbf{Error metric} (\(\downarrow\))} &
\multicolumn{3}{c}{\cellcolor{accbg}\textbf{Accuracy metric} (\(\uparrow\))} \\
 & & &
\textbf{AbsRel} & \textbf{SqRel} & \textbf{RMSE} & \textbf{RMSElog} &
\(\boldsymbol{\delta<1.25}\) & \(\boldsymbol{\delta<1.25^2}\) & \(\boldsymbol{\delta<1.25^3}\) \\
\midrule

ProDepth \cite{woo2024prodepth} &  & low & \cellcolor{lightgreen} 0.086 & \cellcolor{lightgreen} 0.629 & \cellcolor{lightgreen} 4.139 & \cellcolor{lightyellow}  0.166 & \cellcolor{lightyellow} 0.918 & \cellcolor{lightyellow} 0.969 & \cellcolor{lightyellow} 0.984 \\
ManyDepth \cite{watson2021temporal}    &    & low & 0.098 & 0.770 & 4.459 & 0.176 & 0.900 & 0.965 & 0.983 \\
Ours  &    & low &  \cellcolor{lightyellow} 0.089  & \cellcolor{lightyellow}  0.666  &   \cellcolor{lightyellow} 4.192  & \cellcolor{lightgreen}  0.166  &  \cellcolor{lightgreen}  0.918  &   \cellcolor{lightgreen}  0.970  &   \cellcolor{lightgreen}  0.985   \\
Ours  & \checkmark   & low & 0.096  &   0.679  &   4.300  &   0.178  &   0.903  &   0.967  &   0.984   \\
\midrule

ManyDepth \cite{watson2021temporal}   & & medium & 0.093 & 0.715 & 4.245 & 0.172 & 0.909 & 0.966 & 0.983  \\
Ours   &   & medium &  \cellcolor{lightgreen}0.083  &   \cellcolor{lightgreen}0.620  &   \cellcolor{lightgreen}4.020 &   \cellcolor{lightgreen}0.161  &   \cellcolor{lightgreen}0.928  &   \cellcolor{lightgreen}0.972  &   \cellcolor{lightgreen}0.985   \\
Ours   & \checkmark  & medium & 0.091  &   0.634  &   4.132  &   0.173  &   0.914  &   0.970  &   0.984   \\
\bottomrule
\end{tabular}%
}
\end{table*}



\begin{table*}[t]
\centering
\caption{Quantitative comparison with depth foundation models on KITTI.}
\label{tab:foundation}
\small
\setlength{\tabcolsep}{4pt}
\renewcommand{\arraystretch}{1.1}

\begin{tabular}{l|c|c|c c c c|c c c}
\toprule
\textbf{Method} &
\textbf{Metric} &
\textbf{Res.} &
\textbf{AbsRel} & \textbf{SqRel} & \textbf{RMSE} & \textbf{RMSElog} &
\(\boldsymbol{\delta<1.25}\) & \(\boldsymbol{\delta<1.25^2}\) & \(\boldsymbol{\delta<1.25^3}\) \\
\midrule

DepthAnythingV2 (ViT-S) \cite{yang2024depth} & &  518$\times$518  & 0.078 & - & - &  -  & 0.936 & - & -  \\
DepthAnythingV2 (ViT-B) \cite{yang2024depth} & & 518$\times$518 & 0.078 & - & - &  -  & 0.939 & - & -  \\
DepthAnythingV2 (ViT-L) \cite{yang2024depth} & & 518$\times$518  & 0.074 & - & - &  -  & 0.946 & - & -  \\
\rowcolor{lightgreen}
Ours (ViT-B) & & 1246$\times$378 &   0.054  &   0.241  &   2.783  &   0.087  &   0.968  &   0.995  &   0.999  \\
\rowcolor{lightgreen}
Ours (ViT-L) & & 1246$\times$378 &   0.055  &   0.249  &   2.701  &   0.085  &   0.973  &   0.996  &   0.999  \\
\midrule
\rowcolor{lightgreen}
Metric3D V2 (ViT-L, Zero-Shot)\cite{liu2024metricdepth} & \checkmark & 1064$\times$616 & 0.052 & - &  - & - & 0.974 & 0.995 & 0.999 \\

\rowcolor{lightyellow}
Ours (ViT-B) & \checkmark & 1246$\times$378 &   0.063  &   0.264  &   2.870  &   0.094  &   0.963  &   0.995  &   0.999  \\

\rowcolor{lightyellow}
Ours (ViT-L) & \checkmark & 1246$\times$378 &   0.063  &   0.277  &   2.803  &   0.092  &   0.968  &   0.995  &   0.999  \\

\bottomrule
\end{tabular}
\end{table*}






\begin{table*}[t]
\centering
\caption{Robustness against training dataset ratio.}
\label{tab:training_ratio}
% \footnotesize
\setlength{\tabcolsep}{3.5pt}
\renewcommand{\arraystretch}{1.1}
\resizebox{0.7\linewidth}{!}{%
\begin{tabular}{c|c|c c c c|c c c}
\toprule
\textbf{Ratio} &
\textbf{Metric} &
\textbf{AbsRel} & \textbf{SqRel} & \textbf{RMSE} & \textbf{RMSElog} &
\(\boldsymbol{\delta<1.25}\) & \(\boldsymbol{\delta<1.25^2}\) & \(\boldsymbol{\delta<1.25^3}\) \\
\midrule
1\% &   &   0.096  &   0.735  &   4.451  &   0.173  &   0.907  &   0.968  &   0.984  \\
5\% &   &   0.093  &   0.709  &   4.324  &   0.170  &   0.913  &   0.969  &   0.984  \\
10\% &   &   0.091  &   0.692  &   4.262  &   0.168  &   0.915  &   0.970  &   0.985  \\
20\% &   &   0.090  &   0.678  &   4.230  &   0.167  &   0.916  &   0.970  &   0.985  \\
50\% &   &   0.089  &   0.660  &   4.190  &   0.166  &   0.917  &   0.970  &   0.985  \\
\midrule
1\% & \checkmark &   0.101  &   0.750  &   4.553  &   0.185  &   0.892  &   0.963  &   0.982  \\
5\% & \checkmark &   0.097  &   0.709  &   4.379  &   0.180  &   0.898  &   0.966  &   0.983  \\
10\% & \checkmark &   0.098  &   0.699  &   4.354  &   0.180  &   0.899  &   0.966  &   0.983  \\
20\% & \checkmark &   0.097  &   0.688  &   4.333  &   0.179  &   0.902  &   0.967  &   0.984  \\
50\% & \checkmark &   0.096  &   0.674  &   4.304  &   0.178  &   0.903  &   0.967  &   0.984  \\
\bottomrule
\end{tabular}%
% \end{adjustbox}
}
\end{table*}

\textbf{Training Details.}
We train for 5 epochs with batch size 12. The base learning rate is \(1\times10^{-4}\) for all components except the scale--shift head, which uses a larger rate to adapt quickly while keeping the foundation model stable. We apply step decay with factor 0.1 after each epoch. Experiments are conducted at both low resolution (\(182\times630\)) and medium resolution (\(364\times1260\)); unless stated otherwise, results use low resolution. Gradient clipping with maximum norm 20 ensures stability.

\textbf{Evaluation Metrics.}
Following standard practice~\cite{eigen2015predicting, godard2019digging}, we report four error metrics and three accuracy thresholds. \emph{Absolute Relative Error (AbsRel)} measures mean per-pixel depth error normalized by ground-truth depth. \emph{Squared Relative Error (SqRel)} squares the normalized residual, emphasizing outliers. \emph{Root Mean Squared Error (RMSE)} quantifies error magnitude in absolute depth units. \emph{RMSE in log space (RMSElog)} assesses relative depth ordering. For accuracy, we report the percentage of pixels with predicted/ground-truth depth ratio below 1.25, \(1.25^2\), and \(1.25^3\). Together these assess both metric accuracy and structural fidelity.

\begin{figure*}
    \centering
    \includegraphics[width=\linewidth]{imgs/Picture3.png}
    \caption{Visual comparison. Failure regions are highlighted in red and successful predictions in green. Metric3D often underestimates tiny and narrow objects, producing smaller predictions even when the structures are partially preserved. In contrast, our framework better captures fine object boundaries and maintains accurate scale for small objects.}
    \label{fig:viscomp}
\end{figure*}

\subsection{Comparison with Depth Foundation Models}

Since our framework adapts a depth foundation model, Table~\ref{tab:foundation} compares against recent foundation approaches. Using a ViT-B encoder, our non-metric model substantially surpasses DepthAnything\,V2~\cite{yang2024depth}---not only at matched ViT-B capacity but also its larger ViT-L variant (AbsRel \(0.054\) vs.\ \(0.074\), \(\delta{<}1.25\) \(0.968\) vs.\ \(0.946\))---showing that self-supervised adaptation on in-domain video improves accuracy beyond what the frozen foundation model provides, without eroding it. More importantly, our \emph{metric} model---which never sees ground-truth depth---is competitive with the supervised metric foundation model Metric3D\,V2~\cite{liu2024metricdepth} (AbsRel \(0.063\) vs.\ \(0.052\)), narrowing the gap to a model trained on LiDAR with a larger encoder. This is the central result: through scale--shift decoupling and GPS supervision alone, an affine-invariant model is turned into a metric predictor that approaches supervised metric foundation models while retaining its fine-detail fidelity.

\subsection{Comparison with Self-Supervised Methods}

Table~\ref{tab:selfsupervised} compares with state-of-the-art self-supervised methods, ProDepth~\cite{woo2024prodepth} and ManyDepth~\cite{watson2021temporal}, on KITTI. We report both non-metric (median-scaled) and metric (direct output) configurations; although our method natively produces metric depth, we include scale-aligned results for fair comparison with prior work. Notably, both baselines aggregate multiple frames at test time through cost volumes, whereas our model predicts from a \emph{single image}. Despite this, at medium resolution our non-metric variant achieves the best overall accuracy across most metrics, and at low resolution it is on par with ProDepth while preserving superior fine structure, as shown qualitatively. Even our metric configuration---receiving no ground-truth depth---attains competitive accuracy, validating GPS-based scale recovery.



\bigbreak


\begin{table}[h]
\centering
\caption{Generalization study on the cityscapes \cite{cordts2016cityscapes} dataset.}
\label{tab:cityscapes_generalization}
\footnotesize
\setlength{\tabcolsep}{3.5pt}
\renewcommand{\arraystretch}{1.05}

\begin{adjustbox}{max width=\columnwidth}
\begin{tabular}{l|c|c c c c|c c c}
\toprule
\textbf{Method} &
\textbf{Metric} &
% \multicolumn{4}{c|}{\cellcolor{errbg}  &
% \multicolumn{3}{c}{\cellcolor{accbg}  \\
\textbf{AbsRel} & \textbf{SqRel} & \textbf{RMSE} & \textbf{RMSElog} &
\(\boldsymbol{\delta<1.25}\) & \(\boldsymbol{\delta<1.25^2}\) & \(\boldsymbol{\delta<1.25^3}\) \\
\midrule

ManyDepth \cite{watson2021temporal} &  & 0.114 & 1.193 & \textbf{6.223} & 0.170 & 0.875 & 0.967 & 0.989 \\
Ours Zero-shot &   & \textbf{0.112}  &   \textbf{1.067}  &   6.291  &   \textbf{0.165}  &   \textbf{0.884}  &   \textbf{0.976}  &   \textbf{0.992}  \\
\quad w/o scale--shift head &   & 0.114  &   1.075  &   6.329  &   0.167  &   0.881  &   0.975  &   0.992  \\
\midrule
Ours Zero-shot& \checkmark &   \textbf{0.326}  &   \textbf{3.610}  &  \textbf{12.073}  &   \textbf{0.429}  &   0.095  &   \textbf{0.677}  &   \textbf{0.936}  \\
\quad w/o scale--shift head & \checkmark &   0.327  &   3.646  &  12.153  &   0.431  &   0.095  &   0.672  &   0.933  \\
\bottomrule
\end{tabular}
\end{adjustbox}
\end{table}


\subsection{Data Efficiency Analysis}

Table~\ref{tab:training_ratio} evaluates data efficiency by training on progressively smaller subsets of KITTI, from 1\% to 50\%. Accuracy improves monotonically yet only marginally with more data, saturating almost immediately: shrinking the training set 50-fold (from 50\% to 1\%) raises AbsRel by just $0.007$ in the non-metric setting ($0.089\!\rightarrow\!0.096$) and $0.005$ in the metric setting ($0.096\!\rightarrow\!0.101$), while the accuracy thresholds stay essentially unchanged. Training on just 1\% of the data thus recovers nearly all of the accuracy obtained with 50\%. This stems from two factors: the strong pretrained representation of the frozen foundation model provides a robust starting point that requires minimal adaptation, and our parameter-efficient LoRA scheme prevents overfitting while learning effectively from few samples. Practically, this means the method is viable in data-scarce settings where collecting large GPS-annotated video is impractical---adaptation needs only enough data to fix the global scale and lightly tune the backbone, not to learn depth from scratch.



\subsection{Cross-Domain Generalization}

Table~\ref{tab:cityscapes_generalization} evaluates zero-shot transfer: trained on KITTI, tested on Cityscapes without fine-tuning. Despite substantial domain shift---different cameras, lighting, and scene statistics---our non-metric variant (median-scaled) outperforms ManyDepth~\cite{watson2021temporal}---trained directly on Cityscapes---on six of the seven metrics (trailing only on RMSE), confirming that our lightweight adaptation preserves the foundation model's zero-shot generalization rather than overfitting to source-domain appearance. The metric configuration, however, degrades sharply, because the GPS-derived absolute scale learned on KITTI does not transfer to a new domain. This is expected and consistent with our design: metric scale is carried by the low-dimensional affine of Section~3.3 and is inherently dataset-specific, requiring GPS or calibration from the target domain. The strong non-metric transfer shows that the \emph{shape} of our predictions generalizes even when the \emph{scale} does not---a direct empirical signature of scale--shift decoupling.



\subsection{Qualitative Analysis}

Figure~\ref{fig:viscomp} presents qualitative comparisons. Our method produces sharper depth boundaries and better-preserved thin structures (e.g., traffic barriers, poles) than ProDepth and Metric3D\,V2, substantially reducing the over-smoothing that affects methods without strong pretrained priors. Our predictions match the geometric fidelity of DepthAnything\,V2 while adding metric scale, validating that teacher SSI distillation preserves the foundation model's edge-aware detail during adaptation.

\subsection{Ablation Study}

Table~\ref{tab:ablation} isolates the contribution of each component, evaluated in the metric, low-resolution configuration at the final training epoch. Starting from the full model, we remove: (i) the scale--shift head, (ii) teacher SSI distillation, and (iii) high-loss pixel rejection. At convergence the differences are modest but consistent. High-loss pixel rejection has the clearest effect, degrading every metric and most strongly the outlier-sensitive ones (AbsRel \(0.096\!\rightarrow\!0.101\), SqRel \(0.679\!\rightarrow\!0.693\), RMSE \(4.300\!\rightarrow\!4.378\)), confirming that discarding unreliable supervision---typically dynamic objects and occlusion boundaries---stabilizes the recovered metric scale. Removing teacher SSI distillation produces a smaller degradation on these same outlier-sensitive metrics (SqRel \(0.679\!\rightarrow\!0.685\), RMSE \(4.300\!\rightarrow\!4.326\)), indicating that it preserves geometric structure the photometric loss alone tends to erode, while leaving the log and threshold metrics essentially unchanged. Figure~\ref{fig:consistency} makes this qualitative: without the consistency loss, depth boundaries soften and thin structures such as tree trunks and poles are eroded, whereas the full model---anchored to the frozen teacher---keeps them sharp. Disabling the CLS-token scale--shift corrector leaves the converged \emph{in-domain} metric error essentially unchanged (AbsRel \(0.096\!\rightarrow\!0.097\), with a marginal SqRel improvement to \(0.671\)), suggesting that on KITTI the remaining scale pathway already absorbs most of the global affine. Its contribution instead surfaces under domain shift: on Cityscapes (Table~\ref{tab:cityscapes_generalization}) the head improves \emph{every} error metric---and never degrades an accuracy threshold---in both the non-metric and metric settings (non-metric RMSE \(6.329\!\rightarrow\!6.291\), \(\delta{<}1.25\) \(0.881\!\rightarrow\!0.884\); metric SqRel \(3.646\!\rightarrow\!3.610\), RMSE \(12.153\!\rightarrow\!12.073\)), indicating that the explicit corrector generalizes the global affine more robustly than the implicit pathway when the input distribution changes---justifying its inclusion. Overall, the full configuration is best or tied on four of the seven metrics, and the components are complementary, with high-loss rejection contributing most to final metric accuracy.


\begin{table}[t]
\centering
\caption{Ablation study on KITTI Eigen (metric output, low resolution, final epoch). Best per column in \textbf{bold} (ties broken toward the full model).}
\label{tab:ablation}
\footnotesize
\setlength{\tabcolsep}{3.5pt}
\renewcommand{\arraystretch}{1.05}

\begin{adjustbox}{max width=\columnwidth}
\begin{tabular}{l|c c c c|c c c}
\toprule
\textbf{Configuration} &
\textbf{AbsRel} & \textbf{SqRel} & \textbf{RMSE} & \textbf{RMSElog} &
\(\boldsymbol{\delta<1.25}\) & \(\boldsymbol{\delta<1.25^2}\) & \(\boldsymbol{\delta<1.25^3}\) \\
\midrule
\rowcolor{probg}
Full &   \textbf{0.096}  &   0.679  &   \textbf{4.300}  &   0.178  &   0.903  &   \textbf{0.967}  &   \textbf{0.984}  \\
w/o scale--shift head        &   0.097  &   \textbf{0.671}  &   4.313  &   0.179  &   0.901  &   0.967  &   0.984  \\
w/o teacher SSI distillation &   0.096  &   0.685  &   4.326  &   \textbf{0.177}  &   \textbf{0.905}  &   0.967  &   0.984  \\
w/o high-loss rejection      &   0.101  &   0.693  &   4.378  &   0.182  &   0.900  &   0.966  &   0.983  \\
\bottomrule
\end{tabular}
\end{adjustbox}
\end{table}

\begin{figure*}[t]
    \centering
    \includegraphics[width=\linewidth]{imgs/consistency_visual.png}
    \caption{Qualitative effect of the teacher SSI (consistency) distillation loss on KITTI Eigen samples. We compare the full model against the variant trained without the consistency loss (the ``w/o teacher SSI distillation'' row of Table~\ref{tab:ablation}). Removing the loss softens depth boundaries and erodes thin structures---tree trunks, poles, and object silhouettes---which the full model, distilled from the frozen foundation-model teacher, preserves. This mirrors the quantitative degradation on the outlier-sensitive error metrics.}
    \label{fig:consistency}
\end{figure*}


\section{Conclusion}\label{sec6}
This paper introduced a self-supervised framework that turns a frozen affine-invariant depth foundation model into a metric, single-image depth predictor without any depth annotations. We identified the core obstacle---the scale-invariance of self-supervised photometric losses, which both prevents metric scale recovery and erodes the pretrained model's geometric structure---and resolved it through a single principle: decoupling metric scale from depth shape. A \texttt{CLS}-conditioned per-image affine head isolates the two global degrees of freedom (scale and shift) that make a relative prediction metric, GPS supervision injects absolute scale through the pose branch, and teacher scale-shift-invariant distillation protects fine structure from photometric erosion, all realized in under $1\%$ of trainable parameters via LoRA.

Extensive experiments demonstrate that this lightweight adaptation yields metric depth competitive with both self-supervised methods and supervised metric foundation models, while preserving the edge fidelity and cross-domain generalization of the underlying model and remaining effective with as little as $1\%$ of the training data. By showing that metric scale can be grafted onto a relative foundation model through self-supervision alone, our work offers a practical, annotation-free path to metric depth. The main limitation is that GPS-derived scale is dataset-specific and does not transfer zero-shot across domains; future work will investigate scale priors that generalize and the extension of scale--shift decoupling to other geometric foundation tasks.

%\bibliographystyle{cas-model2-names}
\bibliographystyle{IEEEtran}
\bibliography{cas-refs}



%\newpage

%\section{Biography Section}
%If you have an EPS/PDF photo (graphicx package needed), extra braces are
% needed around the contents of the optional argument to biography to prevent
% the LaTeX parser from getting confused when it sees the complicated
% $\backslash${\tt{includegraphics}} command within an optional argument. (You can create
% your own custom macro containing the $\backslash${\tt{includegraphics}} command to make things
% simpler here.)
% 
%\vspace{11pt}
%
%\bf{If you include a photo:}\vspace{-33pt}
%\begin{IEEEbiography}[{\includegraphics[width=1in,height=1.25in,clip,keepaspectratio]{fig1}}]{Michael Shell}
%Use $\backslash${\tt{begin\{IEEEbiography\}}} and then for the 1st argument use $\backslash${\tt{includegraphics}} to declare and link the author photo.
%Use the author name as the 3rd argument followed by the biography text.
%\end{IEEEbiography}
%
%\vspace{11pt}
%
%\bf{If you will not include a photo:}\vspace{-33pt}
%\begin{IEEEbiographynophoto}{John Doe}
%Use $\backslash${\tt{begin\{IEEEbiographynophoto\}}} and the author name as the argument followed by the biography text.
%\end{IEEEbiographynophoto}

\vfill

\end{document}


